\scanpage{49}
\mspage{40}

D'ailleurs les implications
\[
(X_j\subset X_{\deleted{\unsure{j}}i})\Longrightarrow(X_j^{\circ}\subset X_{\deleted{\unsure{j}}i})\overset{\inserted{\text{si }X_j^{\circ}\ne\varnothing}}{\Longrightarrow}(X_j^{\circ}\cap X_{\deleted{\unsure{j}}i})
\]
sont triviales. Donc il reste à prouver, si $X_j^{\circ}\ne\varnothing$
\[
\deleted{\unsure{X_j^{\circ}|X_i^{\circ}}\cap X_j\Longrightarrow i\leq j}\qquad\inserted{X_i^{\circ}}
\]
\[
X_j^{\circ}\cap X_i\ne\varnothing\Longrightarrow j\leq\deleted{\unsure{j}}i.
\]
Mais
\[
X_i=\bigcup_{i'\in I'(i)}X_{i'}^{\circ}\quad\deleted{\text{est saturé pour}}\ \inserted{\text{par }3')},
\]
et comme les $X_{i'}^{\circ}$ $(i'\in I')$ forment une partition de $|\Phi|$, \deleted{\unsure{on en déduit bien}} $X_j^{\circ}$ (pour $j\in I'$ donné) ne peut rencontrer $X_i$ que si $j\deleted{\in I'(i)}\in I'(i)$ i.e. $j\leq i$, qed.

Ainsi on a prouvé
\[
(i)\Longleftrightarrow(iii)\Longleftrightarrow(iv)
\]
en même temps que le corollaire, \unsure{dont} je signale tout de suite le

\uline{Corollaire 2} Supposons \inserted{(dans les conditions équiv. de la prop)} que $\forall i\in I$, on ait $X_i^{\circ}\ne\varnothing$.
Alors l'application
\[
i\longmapsto X_i:\varphi:I\longrightarrow\unsure{\mathscr P|\mathcal P}(X)
\]
est injective, et est un plongement d'ensembles ordonnés : $\forall i,j\in I$, on a
\[
X_i\unsure{\Longleftrightarrow|\subset}X_j\Longleftrightarrow i\leq j.
\]
\inserted{Ainsi,} Les conditions envisagées (i) \deleted{\unsure{de}} \inserted{(plus $\forall i\in I$, $X_i^{\circ}\ne\varnothing$)} sont \deleted{\unsure{clairement}} néc. et suff. pour que \deleted{\unsure{qui définit}} \deleted{soit} $(X_i)_{i\in I}$ \deleted{\unsure{on ait}} soit une \uline{figure ensembliste} indexée par l'ens ord $I$.
\deleted{(\unsure{On a dit (iv)})}

Il reste à voir l'équivalence avec (ii).
% The following final block is struck through horizontally and diagonally.
\deleted{La condition (iii \unsure{b}) \unsure{est} $(x\in|\Phi|)$ dans $3')$.}
\deleted{Le fait que chaque $I^x$ ait un él. \unsure{minimum|minimal} revient à l'\unsure{énoncé} on a}
\[
\deleted{|\Phi|=\bigcup_{i\in I}X_i^{\circ},\quad\text{i.e. }1').}
\]
% archive sheet number: 49
